\documentclass[11pt, a4paper]{article}

% --- Packages ---
\usepackage[utf8]{inputenc}
\usepackage[T1]{fontenc}
\usepackage{geometry}
\geometry{top=1in, bottom=1in, left=1in, right=1in}
\usepackage{xcolor}
\usepackage{listings}
\usepackage{hyperref}
\usepackage{mdframed}
\usepackage{minted}
\usepackage{ifthen}
\usepackage{amsmath}

\newdimen\longline
\longline=\textwidth\advance\longline-6cm
\def\LayoutTextField#1#2{#2}
\def\lbl#1{\hbox to 5cm{#1\hfill\strut}}%
\def\labelline#1#2#3{\lbl{#1}\vbox{\hbox{\TextField[bordercolor=white,name=#2,width=#3]{\null}}\kern2pt\hrule}}
\def\q#1#2{\hbox to \hsize{\labelline{#1}{#2}{\longline}}\vskip1.4ex}



% --- Code Listing Styling ---
\definecolor{codegreen}{rgb}{0,0.6,0}
\definecolor{codegray}{rgb}{0.5,0.5,0.5}
\definecolor{codepurple}{rgb}{0.58,0,0.82}
\definecolor{backcolour}{rgb}{0.95,0.95,0.96}

\lstdefinestyle{mystyle}{
    backgroundcolor=\color{backcolour},   
    commentstyle=\color{codegreen},
    keywordstyle=\color{blue},
    numberstyle=\tiny\color{codegray},
    stringstyle=\color{codepurple},
    basicstyle=\ttfamily\footnotesize,
    breakatwhitespace=false,         
    breaklines=true,                 
    captionpos=b,                    
    keepspaces=true,                 
    numbers=none,                    
    numbersep=5pt,                  
    showspaces=false,                
    showstringspaces=false,
    showtabs=false,                  
    tabsize=4
}
\lstset{style=mystyle}

\title{\textbf{Lab Exercise 1: Measuring the Syscall Tax}}
\author{Gianmaria Del Monte}
\date{}

\begin{document}

\maketitle

\section*{Background}
In the lecture, we discussed the ``Syscall Tax.'' Every time an application invokes a standard \texttt{read()} or \texttt{write()} command, the CPU must transition from \textbf{User Mode} (Ring 3) to \textbf{Kernel Mode} (Ring 0). 

While a single mode switch is incredibly fast, the cumulative cost of millions of switches can destroy performance. In this lab, you will physically quantify this cost and explore how both manual batching and the C Standard Library (\texttt{libc}) work to mitigate this overhead.

\section*{Part 1: Environment Setup}
First, clone the lab repository and navigate to the Exercise 1 directory:

\begin{minted}{shell}
git clone https://github.com/gmgigi96/icscexercises.git
cd icscexercises/exercise1
\end{minted}

Next, compile the programs:
\begin{minted}{shell}
make
\end{minted}

\section*{Part 2: Quantifying the Syscall Tax}
We will use \texttt{strace -c} to intercept system calls and provide a statistical summary.

\subsection*{Task 2.1: Generate the dataset}
Create a \textbf{1 MB} dummy file. While small, reading this 1 byte at a time will trigger over 1 million system calls, which is enough to measure the ``tax'' without taking hours:

\begin{minted}{shell}
dd if=/dev/urandom of=testfile.bin bs=1M count=1
\end{minted}

\subsection*{Task 2.2: The 1-Byte ``Torture Test''}
Run the program reading exactly 1 byte at a time. This forces a mode switch for \textit{every single byte} in the 1 MB file. 

\begin{minted}{shell}
strace -c ./read_test testfile.bin 1
\end{minted}

\textit{Note: This may take 60--90 seconds. Please be patient!}

\subsubsection*{Calculation: How much does a single syscall cost?}
Look at the \texttt{read} row in your \texttt{strace} output. 

\begin{enumerate}
    \item \textbf{Total Time (s):} Record the value from the \texttt{seconds} column: \rule{2cm}{0.4pt}
    \item \textbf{Total Calls:} Record the value from the \texttt{calls} column: \rule{2cm}{0.4pt}
    \item \textbf{Average Latency:}
    \[ \text{Cost per syscall} = \frac{\text{Total seconds}}{\text{Total calls}} = \rule{3cm}{0.4pt} \text{ seconds} \]
\end{enumerate}

\section*{Part 3: Scaling Up and Breaking the Bottleneck}
Now that you have measured the massive cost of a single byte read, we will see how batching allows us to handle much larger datasets.

\subsection*{Task 3.1: Generate a larger dataset}

Create a \textbf{100 MB} file. This is 100x larger than the previous file.
\begin{minted}{shell}
dd if=/dev/urandom of=large_file.bin bs=1M count=100
\end{minted}

\subsection*{Task 3.2: Batching with \texttt{read()}}
Run the test using a 4KB buffer on the large file. Even though the file is 100x larger than the first one, observe the total time:

\begin{minted}{shell}
time ./read_test large_file.bin 4096
\end{minted}

\section*{Part 4: The C Library Rescue}
Programmers often need to process files byte-by-byte (e.g., character parsing). To make this efficient without manual buffer management, \texttt{libc} provides \textbf{Buffered I/O}. 

Look at \texttt{read\_libc.c}. It loops millions of times reading 1 byte, but it uses \texttt{fread()} instead of the raw \texttt{read()} system call. Trace its execution on the large file:

\begin{minted}{shell}
strace -c ./read_libc large_file.bin
\end{minted}

\textit{Observe the number of \texttt{read} calls. Even though the C code asks for 1 byte at a time, \texttt{libc} performs its own batching in user-space. How many syscalls did the kernel actually see?}

\section*{Part 5: Lab Questions}

\noindent \textbf{Question 1:} In Task 2.2 (1-byte read), the CPU was likely pinned at 100\%, but the disk was barely active. What exactly was the CPU doing during those 60--90 seconds?

\vspace{0.2cm}

\noindent \textbf{Question 2:} Based on your \texttt{strace} results, explain the difference between a ``function call'' to the C library (\texttt{fread}) and a ``system call'' to the kernel (\texttt{read}). Why is one so much faster than the other?

\end{document}